% GENERATED FILE. Edit canonical lessons, then run npm run generate:books.
% canonical-source-hash: fnv1a64:6c04b5cc799cbed9
% canonical-lessons: JA-C108-kai, JA-C108-tai, JA-C108-same, JA-C108-risu, JA-C108-inoshishi

\chapter{Shellfish, Sea bream, Shark, Squirrel, Wild boar}
\label{ch:ja-shellfish-sea-bream-shark-squirrel-wild-boar}
\hlchaptermodality{108}

\begin{chapteropening}
\textbf{By the end of this chapter:} \emph{I can say a shellfish, a sea bream, a shark, a squirrel and a wild boar.}

Complete the last lesson of chapter 108: I can say a shellfish, a sea bream, a shark, a squirrel and a wild boar.
\end{chapteropening}

\section[\texorpdfstring{\ja{かい} (kai)}{かい (kai)}]{\ja{かい} (kai) — a shellfish}

\label{lesson:JA-C108-kai}



\begin{quote}
Before the new one: say the Japanese for dew; the rainy season, then the Japanese for a slope.
\end{quote}

\subsection*{You'll want to know: \ja{かい}}
\textbf{\ja{かい}} — \emph{kai} — \textquotedblleft{}a shellfish\textquotedblright{}.

\begin{grammarlens}[title={What you've built}]
One more word for this chapter.
\end{grammarlens}

\subsection*{Your turn}
\begin{itemize}
\raggedright
  \item \emph{Say it:} \emph{kai}
  \item \emph{Say it:} \emph{kai}, once more
  \item \emph{Say it:} say \emph{kai}
  \item \emph{Recall:} say the Japanese for dew; the rainy season, then the Japanese for a slope, then say \emph{kai} again
\end{itemize}

\begin{tcolorbox}[breakable,colback=teal!4,colframe=teal!35!black,title={Before you move on}]
What does \ja{かい} mean? (\textquotedblleft{}A shellfish\textquotedblright{}.) Say it once more.
\end{tcolorbox}
\section[\texorpdfstring{\ja{たい} (tai)}{たい (tai)}]{\ja{たい} (tai) — a sea bream}

\label{lesson:JA-C108-tai}



\begin{quote}
Before the new one: say the Japanese for a slope, then the Japanese for a shellfish.
\end{quote}

\subsection*{You'll want to know: \ja{たい}}
\textbf{\ja{たい}} — \emph{tai} — \textquotedblleft{}a sea bream\textquotedblright{}.

\begin{grammarlens}[title={What you've built}]
One more word for this chapter.
\end{grammarlens}

\subsection*{Your turn}
\begin{itemize}
\raggedright
  \item \emph{Say it:} \emph{tai}
  \item \emph{Say it:} \emph{tai}, once more
  \item \emph{Say it:} say \emph{tai}
  \item \emph{Recall:} say the Japanese for a slope, then the Japanese for a shellfish, then say \emph{tai} again
\end{itemize}

\begin{tcolorbox}[breakable,colback=teal!4,colframe=teal!35!black,title={Before you move on}]
What does \ja{たい} mean? (\textquotedblleft{}A sea bream\textquotedblright{}.) Say it once more.
\end{tcolorbox}
\section[\texorpdfstring{\ja{さめ} (same)}{さめ (same)}]{\ja{さめ} (same) — a shark}

\label{lesson:JA-C108-same}



\begin{quote}
Before the new one: say the Japanese for a shellfish, then the Japanese for a sea bream.
\end{quote}

\subsection*{You'll want to know: \ja{さめ}}
\textbf{\ja{さめ}} — \emph{same} — \textquotedblleft{}a shark\textquotedblright{}.

\begin{grammarlens}[title={What you've built}]
One more word for this chapter.
\end{grammarlens}

\subsection*{Your turn}
\begin{itemize}
\raggedright
  \item \emph{Say it:} \emph{same}
  \item \emph{Say it:} \emph{same}, once more
  \item \emph{Say it:} say \emph{same}
  \item \emph{Recall:} say the Japanese for a shellfish, then the Japanese for a sea bream, then say \emph{same} again
\end{itemize}

\begin{tcolorbox}[breakable,colback=teal!4,colframe=teal!35!black,title={Before you move on}]
What does \ja{さめ} mean? (\textquotedblleft{}A shark\textquotedblright{}.) Say it once more.
\end{tcolorbox}
\section[\texorpdfstring{\ja{りす} (risu)}{りす (risu)}]{\ja{りす} (risu) — a squirrel}

\label{lesson:JA-C108-risu}



\begin{quote}
Before the new one: say the Japanese for a sea bream, then the Japanese for a shark.
\end{quote}

\subsection*{You'll want to know: \ja{りす}}
\textbf{\ja{りす}} — \emph{risu} — \textquotedblleft{}a squirrel\textquotedblright{}.

\begin{grammarlens}[title={What you've built}]
One more word for this chapter.
\end{grammarlens}

\subsection*{Your turn}
\begin{itemize}
\raggedright
  \item \emph{Say it:} \emph{risu}
  \item \emph{Say it:} \emph{risu}, once more
  \item \emph{Say it:} say \emph{risu}
  \item \emph{Recall:} say the Japanese for a sea bream, then the Japanese for a shark, then say \emph{risu} again
\end{itemize}

\begin{tcolorbox}[breakable,colback=teal!4,colframe=teal!35!black,title={Before you move on}]
What does \ja{りす} mean? (\textquotedblleft{}A squirrel\textquotedblright{}.) Say it once more.
\end{tcolorbox}
\section[\texorpdfstring{\ja{いのしし} (inoshishi)}{いのしし (inoshishi)}]{\ja{いのしし} (inoshishi) — a wild boar}

\label{lesson:JA-C108-inoshishi}



\begin{quote}
Before the new one: say the Japanese for a shark, then the Japanese for a squirrel.
\end{quote}

\subsection*{You'll want to know: \ja{いのしし}}
\textbf{\ja{いのしし}} — \emph{inoshishi} — \textquotedblleft{}a wild boar\textquotedblright{}.

\begin{grammarlens}[title={What you've built}]
One more word for this chapter.
\end{grammarlens}

\subsection*{Your turn}
\begin{itemize}
\raggedright
  \item \emph{Say it:} \emph{inoshishi}
  \item \emph{Say it:} \emph{inoshishi}, once more
  \item \emph{Say it:} say \emph{inoshishi}
  \item \emph{Recall:} say the Japanese for a shark, then the Japanese for a squirrel, then say \emph{inoshishi} again
\end{itemize}

\begin{tcolorbox}[breakable,colback=teal!4,colframe=teal!35!black,title={Before you move on}]
What does \ja{いのしし} mean? (\textquotedblleft{}A wild boar\textquotedblright{}.) Say it once more.
\end{tcolorbox}
